% ============================== Part E ==============================
\problem{E1}{When the data can't determine the answer}

\begin{given}
A linear model $\vd = A\mathbf{s} + \mathbf n$, with unknowns $\mathbf s$ and noise covariance $N$. Weighted least squares
gives $\hat{\mathbf s} = [A^\T N^{-1}A]^{-1}A^\T N^{-1}\vd$, which needs $A^\T N^{-1}A$ to be invertible.
\source{\cHBLwls.}
\end{given}

\need{A3 (the matrix least-squares solution).}

\begin{enumerate}
  \item One datum measures the sum of two unknowns: $d = s_1 + s_2 + n$, so $A = (1\ \ 1)$. Show that $A^\T A$ is
    singular. Describe all the solutions that fit the datum perfectly.
  \item Two data, $A = \mathrm{diag}(1, 0.01)$, $N = \sigma^2 I$. Show that $\hat s_2$ amplifies the noise 100-fold.
  \item Our $\kappa$ map has 144 unknown pixels and about 150 galaxies, each with noise $\approx 0.05$ against a subhalo
    signal of $\approx 0.01$. Which of (a) and (b) describes it?
\end{enumerate}

\solution{
(a) $A^\T A = \begin{pmatrix}1&1\\1&1\end{pmatrix}$ has determinant 0. Any $(s_1, s_2)$ with $s_1 + s_2 = d$ fits; the
data fix one combination and say nothing about $s_1 - s_2$.
(b) $\hat s_2 = d_2/0.01 = s_2 + 100\,n_2$: a direction the data barely respond to is recovered with 100 times the noise.
(c) Both. Small-scale patterns of $\kappa$ produce shears that nearly cancel at the galaxies (small singular values of
$A$, as in (b)), and the noisy, nearly degenerate system has many near-equally good solutions (as in (a)). Something
other than the data must choose between them: a prior.
}

% --------------------------------------------------------------------
\problem{E2}{A Gaussian prior makes the problem well posed}

\begin{given}
Add a Gaussian prior $\mathbf s \sim \mathcal N(0, S)$ to the model of E1. For a Gaussian signal and Gaussian noise, the
linear minimum-variance (Wiener) filter is $W = SA^\T(ASA^\T + N)^{-1}$, and it is also the maximum and the mean of the
Gaussian posterior.
\source{\cZaroubiWF; \cJeffreyWF.}
\end{given}

\need{A3 (completing the square in matrix form), E1.}

\begin{enumerate}
  \item Write $-2\ln p(\mathbf s\mid\vd) = (\vd - A\mathbf s)^\T N^{-1}(\vd - A\mathbf s) + \mathbf s^\T S^{-1}\mathbf s + \text{const}$.
    Complete the square to show that the posterior is Gaussian with precision
    $\mathsf H = S^{-1} + A^\T N^{-1}A$ and mean $\hat{\mathbf s} = \mathsf H^{-1}A^\T N^{-1}\vd$.
  \item Show that $\hat{\mathbf s} = W\vd$, using the identity
    $(S^{-1} + A^\T N^{-1}A)^{-1}A^\T N^{-1} = SA^\T(ASA^\T + N)^{-1}$. Verify the identity by multiplying both sides on
    the left by $S^{-1} + A^\T N^{-1}A$ and on the right by $ASA^\T + N$.
  \item Redo E1(a) with $S = \sigma_s^2 I$ and $N = \sigma_n^2$. Is $\mathsf H$ invertible now? What is $\hat{\mathbf s}$?
\end{enumerate}

\solution{
(a) The $\mathbf s$-dependent terms are $\mathbf s^\T(A^\T N^{-1}A + S^{-1})\mathbf s - 2\mathbf s^\T A^\T N^{-1}\vd$,
which equals $(\mathbf s - \hat{\mathbf s})^\T\mathsf H(\mathbf s - \hat{\mathbf s})$ + const with
$\hat{\mathbf s} = \mathsf H^{-1}A^\T N^{-1}\vd$. A Gaussian's mean and maximum coincide.
(b) Left side times the two factors: $A^\T N^{-1}(ASA^\T + N) = A^\T N^{-1}ASA^\T + A^\T$. Right side:
$(S^{-1} + A^\T N^{-1}A)SA^\T = A^\T + A^\T N^{-1}ASA^\T$. They agree.
(c) $\mathsf H = \sigma_s^{-2}I + \sigma_n^{-2}\begin{pmatrix}1&1\\1&1\end{pmatrix}$ has eigenvalues
$\sigma_s^{-2} + 2\sigma_n^{-2}$ and $\sigma_s^{-2}$, both positive, so it is invertible. By symmetry
$\hat s_1 = \hat s_2 = \frac{d\,\sigma_s^2}{2\sigma_s^2 + \sigma_n^2}$: the datum is split equally, shrunk toward zero,
and the undetermined direction $s_1 - s_2$ is set to its prior value, 0.
}

% --------------------------------------------------------------------
\problem{E3}{One mode at a time: what the Wiener filter keeps}

\begin{given}
If $S$ and $N$ are diagonal in the same basis, for example Fourier modes of a stationary field, the filter acts on
each mode separately and suppresses the noise by the ratio prior/(prior + noise). It is the linear estimate with the
smallest residual variance.
\source{\cZaroubiMSE.}
\end{given}

\need{E2.}

\begin{enumerate}
  \item One mode, $d = s + n$, with $s \sim \mathcal N(0, S)$ and $n \sim \mathcal N(0, N)$. For an estimate
    $\hat s = w\,d$, show that the mean squared error $\E[(\hat s - s)^2] = (1-w)^2S + w^2N$, and that it is minimised
    at $w = S/(S+N)$.
  \item For a \emph{fixed} true $s$, show that $\E[\hat s\mid s] = ws$. In what sense is the Wiener estimate biased?
  \item Large scales have $S \gg N$ and small scales $S \ll N$. Describe the filtered map.
  \item The subhalo in our likelihood map recovers about 40\% of its true peak value. What value of $w$ does that
    suggest for the scales it occupies?
\end{enumerate}

\solution{
(a) $\hat s - s = (w-1)s + wn$; with $s$ and $n$ independent, the variance is $(1-w)^2S + w^2N$. Setting the
derivative to zero, $-2(1-w)S + 2wN = 0$, gives $w = S/(S+N)$.
(b) $\E[w(s+n)\mid s] = ws$: every mode is shrunk toward zero by $w < 1$. The estimate is biased for any given field,
but over the ensemble of fields its error is minimal. This is the trade made on purpose.
(c) Large scales pass almost unchanged and small scales are suppressed, so the map is a smoothed, slightly shrunk
version of the truth.
(d) $w \approx 0.4$: at the subhalo's scale the signal and noise powers are comparable ($S/N \approx 0.7$).
}

% --------------------------------------------------------------------
\problem{E4}{A smoothness penalty is a power spectrum}

\begin{given}
Our map prior is
$-\ln p(\kappa) = \frac{1}{2\sk^2}\big[\sum_j\kappa_j^2 + \ell^2\sum_{\langle jj'\rangle}(\kappa_j - \kappa_{j'})^2\big]$,
summed over pixels $j$ and neighbouring pairs $\langle jj'\rangle$. \emph{Derived here}; it is the regulariser of
\texttt{likelihood\_kappa} in \texttt{h\_cluster\_analysis.py}.
\end{given}

\need{E2, E3 (a Gaussian prior with covariance $S$; per-mode filtering).}

\begin{enumerate}
  \item On a periodic 1D chain of $M$ pixels, insert a Fourier mode $\kappa_j = a\,e^{ikj}$. Show that
    $\sum_j |\kappa_{j+1} - \kappa_j|^2 = 4\sin^2(k/2)\sum_j|\kappa_j|^2$.
  \item Conclude that each mode has prior variance $P(k) = \sk^2/[1 + 4\ell^2\sin^2(k/2)] \approx \sk^2/(1+\ell^2k^2)$
    for small $k$.
  \item What do $\sk$ and $\ell$ each control? Which scales does the prior suppress?
\end{enumerate}

\solution{
(a) $\kappa_{j+1} - \kappa_j = a\,e^{ikj}(e^{ik} - 1)$ and $|e^{ik} - 1|^2 = 2 - 2\cos k = 4\sin^2(k/2)$.
(b) Modes are independent (they diagonalise the quadratic form), and mode $k$ has precision
$[1 + 4\ell^2\sin^2(k/2)]/\sk^2$, whose inverse is the prior variance. For small $k$, $4\sin^2(k/2) \approx k^2$.
(c) $\sk$ sets the overall amplitude the prior allows. $\ell$ sets the scale below which fluctuations are suppressed:
the variance falls by half at $k \approx 1/\ell$. Combined with E3, the Wiener weight per mode is
$P(k)/[P(k) + N(k)]$, which removes small scales first.
}

% --------------------------------------------------------------------
\problem{E5}{\py\ A one-dimensional Wiener filter}

\begin{given}
A toy 1D ``$\kappa$'' made of a broad bump (the main halo) plus a narrow bump (the subhalo) is sampled by noisy point
measurements at random positions, each with noise $\sigma = 0.05$ of a quantity linear in $\kappa$. The MAP map under
the prior of E4 is $\hat\kappa = \mathsf H^{-1}A^\T N^{-1}\vd$ (E2). \emph{Derived here}.
\end{given}

\need{E2, E4.}

\begin{enumerate}
  \item \py\ Build $A$ (a measurement samples the pixel it falls in) and the prior precision of E4. Compute the MAP
    map for $\sk \in \{0.03, 0.1, 0.3\}$ and $\ell \in \{0.5, 1, 2\}$ pixels.
  \item \py\ Which settings recover the narrow bump best? What happens to it for small $\sk$, and for large $\ell$?
  \item \py\ Repeat with 20 noise realizations. For one setting, plot the mean map against the truth. Is the mean map
    equal to the truth?
\end{enumerate}

\solution{
(a)--(b) See the solutions notebook. Small $\sk$ shrinks everything (low peaks); large $\ell$ blurs the narrow bump into
the broad one. Large $\sk$ with small $\ell$ keeps the bump's height but lets the noise through. The best setting
balances the two, just as the tuning did on the real map.
(c) No: the mean map is a smoothed, shrunk version of the truth (E3(b)). Averaging over noise removes the noise, not
the prior's bias. Part F makes this precise with the resolution matrix.
}
